\section{Related Work}

\subsection{Heterogeneous Graph Neural Networks}

Heterogeneous graph neural networks extend graph representation learning to graphs with multiple node and relation types. HAN \cite{wang2019han} uses node-level attention and semantic-level attention to aggregate information from meta-path-based neighborhoods. HGT \cite{hu2020hgt} introduces type-specific projections and attention mechanisms for heterogeneous graph transformers. MAGNN \cite{fu2020magnn} models intra-meta-path and inter-meta-path aggregation to capture fine-grained semantic dependencies. HeCo \cite{wang2021heco} further uses contrastive learning between network schema and meta-path views to improve heterogeneous graph representation learning.

These methods show the value of semantic aggregation, but they are usually designed for clean graphs. Under adversarial structural perturbations, their reliance on observed edges and meta-path-induced neighborhoods may cause corrupted semantics to be propagated. DVCL keeps the semantic modeling ability of heterogeneous graph learning while explicitly introducing topology purification and a feature-induced view for robustness.

\subsection{Robust Graph Neural Networks}

Robust graph neural networks aim to reduce the sensitivity of graph models to adversarial or noisy structures. GCN-Jaccard \cite{wu2019gcnjaccard} removes edges between nodes with low feature similarity. RGCN \cite{zhu2019rgcn} uses Gaussian distributions to model hidden representations and reduce adversarial effects. Pro-GNN \cite{jin2020prognn} jointly learns graph structure and GNN parameters with sparsity and low-rank constraints. GNNGuard \cite{zhang2020gnnguard} assigns edge weights according to feature similarity, and SimP-GCN \cite{jin2021simpgcn} incorporates similarity-preserving constraints into graph convolution.

Most robust GNNs are developed for homogeneous graphs. Directly applying them to heterogeneous graphs may ignore typed relations and semantic paths. DVCL differs by operating on meta-path-induced semantic structures and by using feature-induced evidence as an auxiliary view rather than as a simple edge-pruning rule.

\subsection{Robust Heterogeneous Graph Learning}

Robust heterogeneous graph learning has recently attracted attention because structural attacks can be amplified through heterogeneous semantics. Methods such as RoHe \cite{rohe}, HeteGuard \cite{heteguard}, and HSeCo \cite{hseco} attempt to defend HGNNs by modeling relation-aware perturbations, filtering unreliable semantic connections, or improving semantic topology learning.

These methods are closely related to DVCL, but they still mainly rely on the attacked heterogeneous topology when constructing candidate neighborhoods. DVCL addresses this limitation by adding a feature-induced graph that does not directly depend on perturbed edges. The cross-view objective then encourages topology-view representations and feature-view representations to agree, improving robustness when either view is noisy.

\subsection{Graph Contrastive Learning}

Graph contrastive learning learns representations by maximizing agreement between related graph views. DGI \cite{velickovic2019dgi} contrasts node representations with graph-level summaries. MVGRL \cite{hassani2020mvgrl} contrasts diffusion and adjacency views. GRACE \cite{zhu2020grace} and GCA \cite{zhu2021gca} generate graph augmentations and align node representations across augmented views. In heterogeneous graphs, HeCo \cite{wang2021heco} contrasts schema and meta-path views to learn semantic representations.

DVCL is also based on cross-view agreement, but its views are designed for robustness rather than only self-supervised representation learning. The topology view captures purified heterogeneous semantics, while the feature-induced view provides an attack-resistant reference. The contrastive objective is therefore used to improve consistency between semantic topology and feature similarity under structural attacks.

% ========== 6 Conclusion ==========

